%%%%%%%%%%%%%%%%%%%%%%%%%%%%%%%%%%%%%%%%%%%%%%%%%%%%%%%%%%%%%%%%%%%%%
%% Table S4: TD-DFT Benchmarking
%%%%%%%%%%%%%%%%%%%%%%%%%%%%%%%%%%%%%%%%%%%%%%%%%%%%%%%%%%%%%%%%%%%%%

\begin{table}[H]
\centering
\caption{Benchmarking of xTB-based methods against TD-DFT calculations. 
Vertical excitation energies and singlet-triplet gaps were calculated at the 
$\omega$B97X-D/def2-TZVP level of theory using ORCA 6.1.0, with CPCM solvation in toluene. 
The benchmark set comprises 8 representative TADF emitters spanning different 
donor-acceptor architectures. All structures were optimized at the GFN2-xTB level 
in toluene solvent. Both sTDA-xTB and sTD-DFT-xTB systematically underestimate 
\deltaest\, by approximately 0.77--0.79 eV compared to TD-DFT, but preserve the relative 
ranking of molecules.}
\label{tab:tddft_benchmark}
\footnotesize
\begin{tabular}{l S[table-format=1.3] S[table-format=1.3] S[table-format=1.3] S[table-format=1.3] S[table-format=1.3] S[table-format=+1.3] S[table-format=+1.3]}
\toprule
\textbf{Molecule} & 
{\textbf{TD-DFT}} & 
{\textbf{TD-DFT}} & 
{\textbf{TD-DFT}} & 
{\textbf{sTDA-xTB}} & 
{\textbf{sTD-DFT-xTB}} & 
{\textbf{Deviation}} & 
{\textbf{Deviation}} \\
& 
{$S_1$ [eV]} & 
{$T_1$ [eV]} & 
{$\Delta E_{\text{ST}}$ [eV]} & 
{$\Delta E_{\text{ST}}$ [eV]} & 
{$\Delta E_{\text{ST}}$ [eV]} & 
{sTDA [eV]} & 
{sTD-DFT [eV]} \\
\midrule
2TCz-DPS   & 4.251 & 3.159 & 1.092 & 0.321 & 0.296 & +0.771 & +0.796 \\
4CzIPN     & 3.685 & 2.801 & 0.884 & 0.234 & 0.208 & +0.650 & +0.676 \\
CzS2       & 4.455 & 3.196 & 1.259 & 0.381 & 0.352 & +0.878 & +0.907 \\
DMAC-DPS   & 4.168 & 3.329 & 0.839 & 0.234 & 0.234 & +0.605 & +0.605 \\
DMAC-TRZ   & 4.095 & 3.071 & 1.024 & 0.049 & 0.051 & +0.975 & +0.973 \\
PSPCz      & 4.481 & 3.197 & 1.284 & 0.404 & 0.363 & +0.880 & +0.921 \\
Px2BP      & 3.827 & 2.991 & 0.836 & 0.226 & 0.214 & +0.610 & +0.622 \\
TDBA-DI    & 4.052 & 2.889 & 1.163 & 0.380 & 0.344 & +0.783 & +0.819 \\
\midrule
\textbf{Mean}     & 4.127 & 3.079 & 1.048 & 0.279 & 0.258 & +0.769 & +0.790 \\
\textbf{MAE}      & \multicolumn{3}{c}{---} & \multicolumn{2}{c}{0.769} & \multicolumn{2}{c}{0.790} \\
\textbf{RMSE}     & \multicolumn{3}{c}{---} & \multicolumn{2}{c}{0.780} & \multicolumn{2}{c}{0.801} \\
\bottomrule
\end{tabular}
\end{table}

\textbf{Notes:}
\begin{itemize}
\item \textbf{Computational details:} TD-DFT calculations performed with ORCA 6.1.0 using 
$\omega$B97X-D functional with def2-TZVP basis set and CPCM solvation model (toluene, $\varepsilon=2.374$). 
Vertical excitations calculated on GFN2-xTB/GBSA(toluene) optimized geometries.

\item \textbf{Systematic underestimation:} Both sTDA-xTB and sTD-DFT-xTB consistently 
underestimate \deltaest\ compared to TD-DFT (mean deviation $\sim\qty{0.77}{\electronvolt}$--$\qty{0.79}{\electronvolt}$). This systematic 
bias does not affect relative ranking for high-throughput screening applications.

\item \textbf{Correlation:} Despite the systematic offset, xTB methods correctly identify 
molecules with smaller vs. larger \deltaest\ values (e.g., Px2BP and 4CzIPN have the smallest 
gaps in both TD-DFT and xTB calculations).

\item \textbf{Computational cost:} TD-DFT calculations require $\sim2-6$ hours per molecule 
on 8 CPU cores, while xTB methods complete in $\sim2-5$ minutes, enabling the screening of 
747 molecules presented in this work.

\item \textbf{Recommendation:} Use xTB methods for initial high-throughput screening to 
identify promising candidates, followed by TD-DFT validation for quantitative accuracy.
\end{itemize}

